% ==============================================================================
% SECTION 5: CONCLUSION
% ==============================================================================
\section{Conclusion}
\begin{frame}{Synthèse de la 1ère Partie}
    \begin{block}{Ce qu'il faut retenir :}
        \begin{itemize}
            \item \textbf{Prérogatives :} PA60 (avec DP) et PA40 (sans DP). Vous devez connaître vos limites réglementaires.
            \item \textbf{Gestion d'air :} L'autonomie diminue drastiquement avec la profondeur. La \textbf{règle des tiers} est votre assurance vie.
            \item \textbf{Barotraumatismes :} La variation de pression est rapide près de la surface. Vigilance à la descente comme à la remontée.
            \item \textbf{ADD :} Le respect des procédures ne suffit pas toujours. Il faut combattre les facteurs favorisants (froid, effort, profil).
        \end{itemize}
    \end{block}

    \begin{center}
        \large \textbf{Le Niveau 3 est le niveau de l'autonomie complète.} \\
        \vspace{0.5em}
        \textit{Votre sécurité dépend de votre planification et de votre capacité à réagir.}
    \end{center}
\end{frame}

\begin{frame}{Vers la 2ème Partie...}
    Nous avons vu les risques et la théorie.
    
    Dans la seconde partie, nous aborderons la pratique de l'autonomie :
    \begin{itemize}
        \item Comment utiliser son ordinateur et gérer la décompression ?
        \item Comment s'organiser matériellement (bateau, sécu) ?
        \item Comment planifier concrètement la plongée ?
    \end{itemize}
\end{frame}

\begin{frame}
    \centering
    \Huge \textbf{Des questions ?}
\end{frame}